\chapter{Templates and batch processing}
\label{ch:batch}

If you have twenty proteins to prepare the same way, you do not have to answer
the same questions twenty times. Prepare one protein by hand, mark the answers
that change from protein to protein, and let ProPrep work through a list.

A recipe with blanks in it is called a \term{template}. Running it over a list
is called \term{batch processing}. This chapter follows the whole route, using
the example that worked: download a structure, remove its crystal waters, and
repeat that for three proteins.

\begin{glance}
\glanceitem{What it does}{Turns one recorded preparation into a recipe and runs
it on a list of proteins, unattended.}
\glanceitem{You need it when}{You have several similar proteins to prepare the
same way.}
\glanceitem{Before you start}{One finished preparation, recorded as a session
(Chapter~\ref{ch:sessions}). Choose proteins that will be asked the same
questions.}
\glanceitem{Where to find it}{Only from the Terminal:
\texttt{--create-template}, \texttt{--validate-template},
\texttt{--batch-replay}.}
\glanceitem{You get}{A template file, a list file, one folder per protein, and
a report naming every success and failure.}
\end{glance}

\section{Step by step}

\subsection{Step 1: prepare one protein by hand}

Work through your preparation once, normally. ProPrep records it without being
asked (Chapter~\ref{ch:sessions}). Note the line it prints at the start:

\begin{screen}
Session recording to: ./proprep_session_20261121_122825.json
\end{screen}

Our example session had 19 answers: load \pdb{1UBQ} from the PDB, run the PDB
Filter, keep the protein, discard the water, exit.

\begin{tip}
Keep this first preparation as short and plain as possible. Every extra
question is one more thing that can differ for the next protein.
\end{tip}

\subsection{Step 2: make the template}

Copy the session file somewhere tidy if you like, then:

\begin{terminal}
proprep --create-template proprep_session_20261121_122825.json
\end{terminal}

You get the Template Creator, which is the Session Editor
(Section~\ref{sec:sessions-editor}) with one extra command, \typed{v}, for
marking an answer as a blank.

\begin{screen}
╭──────────────────── Template Creator ─────────────────────╮
│ File: ubq_session.json  |  19 interactions  |  no changes │
╰───────────────────────────────────────────────────────────╯
┏━━━━┳━━━━━━━━┳━━━━━━━━━━━━━━━━━━━━━━━━┳━━━━━━━━━━━━━━━━━━━━━━━┳━━━━━━━━━━━━━┓
┃  # ┃ Status ┃ Module                 ┃ Question              ┃ Answer      ┃
┡━━━━╇━━━━━━━━╇━━━━━━━━━━━━━━━━━━━━━━━━╇━━━━━━━━━━━━━━━━━━━━━━━╇━━━━━━━━━━━━━┩
│  5 │   ·    │ Structure Loader - RCSB│ Enter PDB ID to ...   │ 1UBQ        │
│ 16 │   ·    │ PDB Filter             │ Chain A - Water ...   │ Discard ... │
└────┴────────┴────────────────────────┴───────────────────────┴─────────────┘
[e]dit #  [d]elete #|#-#|#-end  [u]ndo #  [v]ariable #  [n]ext  [b]ack
[j]ump #  [p]review  save  quit  help
\end{screen}

\begin{watchout}
Nothing is marked for you. In version 1.23.0 the PDB code is not detected
automatically, so a template with no \typed{v} command in it has no blanks at
all. Mark them yourself.
\end{watchout}

Mark the PDB code, which is answer 5 in our session:

\begin{screen}
> : §ans«v 5»
Marking #5 as template variable
Interaction: Enter PDB ID to download
Original value: 1UBQ (1UBQ)
Variable name (input_protein):
Description (Enter PDB ID to download):
Required? [y/n] (y): §ans«y»
Marked as variable: {{ input_protein }}
\end{screen}

Press \key{Enter} to accept the suggested name \typed{input\_protein}; the
plain-text list format in Section~\ref{sec:batch-list} fills that name
automatically. \prompt{Required} means every run must supply a value.

You can mark more than one answer. We also marked the water question, naming
it ourselves:

\begin{screen}
> : §ans«v 16»
Marking #16 as template variable
Interaction: Chain A - Water retention option
This interaction had these options:
  r: Retain entire component
  ...
  d: Discard entire component <-- original
Variable name (var_16): §ans«water_option»
Description (Chain A - Water retention option):
§ans«Keep (r) or discard (d) crystal waters»
Required? [y/n] (y): §ans«y»
\end{screen}

Type \typed{p} to check your marks, then \typed{save}:

\begin{screen}
Template saved: ubq_session_template.json
2 variable(s), 0 edit(s), 0 deletion(s)

Template variables:
  *input_protein: Enter PDB ID to download (example: 1UBQ)
  *water_option: Keep (r) or discard (d) crystal waters (example: d)
\end{screen}

\begin{watchout}
A template needs at least one blank. Typing \typed{save} with none gives
\prompt{No variables marked. A template needs at least one variable.} and then
offers to save it as an ordinary edited session instead.
\end{watchout}

\begin{watchout}
After \typed{e 5} (edit), the next thing you type is taken as the new answer.
Typing \typed{quit} there stores the word ``quit'' as the PDB code. Finish the
edit first, then quit.
\end{watchout}

\subsection{Step 3: make the list of proteins}

ProPrep offers to build the list for you straight away. Answer \key{y}:

\begin{screen}
Create a batch input file now? [y/n] (y): §ans«y»

Enter variable values for each run. Press Enter with no input to finish.

Run 1:
  input_protein (1UBQ):
  water_option options:
    r: Retain entire component
    s: Select specific residues
    d: Discard entire component
  water_option (d): §ans«d»

Run 2:
  input_protein (): §ans«1AKI»
  water_option (): §ans«r»
...
Run 4:
  input_protein ():

Batch file saved: ubq_session_batch.csv (3 runs)
Run with: proprep --batch-replay ubq_session_template.json --batch-list
ubq_session_batch.csv
\end{screen}

Press \key{Enter} on an empty protein to finish. The file it writes is a plain
table you can also make yourself in a spreadsheet or a text editor:

\begin{screen}
input_protein,water_option
1UBQ,d
1AKI,r
1CRN,d
\end{screen}

If you answer \key{n}, ProPrep tells you the format for later:
\prompt{CSV with headers matching variable names}.

\begin{tip}
If your template has only one blank, and it is named
\typed{input\_protein}, the list can be a plain text file with one PDB code per
line. Lines starting with \texttt{\#} are ignored.
\end{tip}

\subsection{Step 4: make the output folder}

\begin{terminal}
mkdir batch_out
\end{terminal}

\begin{watchout}
Do this first. In version 1.23.0, ProPrep does not create the output folder.
If it does not exist, every run fails at once with
\prompt{Setup error: [Errno 2] No such file or directory: 'batch\_out/run\_01'}
and the batch ends with a Python error report instead of a summary.
\end{watchout}

\subsection{Step 5: check the template against the list}

This takes a second and saves a wasted batch:

\begin{terminal}
proprep --validate-template ubq_session_template.json \
        --batch-list ubq_session_batch.csv
\end{terminal}

\begin{screen}
✓ Template is valid: ubq_session_template.json
\end{screen}

If a blank has no value, ProPrep names it and the run it is missing from, and
stops with an error. A column in your list that matches no blank is only a
warning: \prompt{! Input column 'extra' matches no template variable and will
be ignored (check the header for a typo)}.

\begin{watchout}
Validation checks that every blank has a value. It does \textbf{not} check
that the values make sense. A water option of \typed{x}, which is not one of
\texttt{r}, \texttt{s} or \texttt{d}, passes validation, and then during the
run ProPrep accepts it silently and keeps the waters. Check your list by eye.
\end{watchout}

\subsection{Step 6: run the batch}

\begin{terminal}
proprep --batch-replay ubq_session_template.json \
        --batch-list ubq_session_batch.csv \
        --output-dir batch_out --batch-continue
\end{terminal}

\begin{screen}
═══ Batch Processing Mode ═══

Template: ubq_session_template.json
Runs to process: 3
Template variables:
  - input_protein: Enter PDB ID to download
  - water_option: Keep (r) or discard (d) crystal waters

Replay mode: plain terminal
\end{screen}

A progress bar names the protein being worked on. At the end comes the
summary:

\begin{screen}
═══ Batch Processing Summary ═══

Planned: 3
Successful: 2 (66.7%)
Failed: 1 (33.3%)

┏━━━━━━━━━┳━━━━━━━━━━━┳━━━━━━━━━━━━━━━━━━━━━━━━━━━━━━━━━━━━━━━━━━━━━━━━━━━━━┓
┃ Protein ┃ Status    ┃ Output Directory                                    ┃
┡━━━━━━━━━╇━━━━━━━━━━━╇━━━━━━━━━━━━━━━━━━━━━━━━━━━━━━━━━━━━━━━━━━━━━━━━━━━━━┩
│ 1CRN    │ ✗ Failed  │ batch_cont/run_01                                   │
│         │           │ Replay diverged: no recorded answer for prompt      │
│         │           │ 'Enter your choice'. At this point the recording    │
│         │           │ expects a prompt here: 'Choose option'.             │
│ 1UBQ    │ ✓ Success │ batch_cont/run_02                                   │
│ 1AKI    │ ✓ Success │ batch_cont/run_03                                   │
└─────────┴───────────┴─────────────────────────────────────────────────────┘

Batch report saved to: batch_cont/batch_report_20261121_125042.json
Runs still owed (1): batch_cont/batch_retry_20261121_125042.csv
Re-run them with: proprep --batch-replay ubq_session_template.json
--batch-list batch_cont/batch_retry_20261121_125042.csv
\end{screen}

Three proteins took about one second in total, because downloading and
filtering are quick. The whole batch returns a failure code if any protein
failed, which matters if you call it from a script.

\subsection{Step 7: read the results}

Each protein gets a numbered folder:

\begin{tabularx}{\linewidth}{@{}lL@{}}
\toprule
\textbf{File} & \textbf{What it holds} \\
\midrule
\file{run_01/run.log} & Everything ProPrep printed for that protein, without
  colors. Read this first. \\
\file{run_01/error.log} & Only written when the run failed: the reason and a
  Python error report. \\
\file{run_01/}\dots & The output files, exactly as a run by hand would have
  made them (here \file{1UBQ.pdb} and \file{filtered_structure.pdb}). \\
\file{batch_report_*.json} & The whole batch: the values used for each run and
  whether it worked. \\
\file{batch_retry_*.csv} & Only the proteins that failed or were never
  attempted, ready to use as a new list. \\
\bottomrule
\end{tabularx}

\medskip
To retry the failures, pass the retry file as the list. ProPrep prints that
command for you, as in the summary above.

\section{Choosing proteins that will work}
\label{sec:batch-choose}

A batch replays your \emph{answers}. It works only when the new protein is
asked the \emph{same questions} as the recorded one.

This is why \pdb{1CRN} failed above. Crambin has no crystal waters, so the
PDB Filter never asks about water. The recording has an answer for a question
that is not asked, and so the answers slip out of step. ProPrep notices and
stops that run:

\begin{screen}
Replay diverged: no recorded answer for prompt 'Enter your choice'. At this
point the recording expects a prompt here: 'Choose option'.
\end{screen}

\prompt{Replay diverged} always means the same thing: ProPrep was asked
something the recording did not expect. Read it as ``this protein is not like
the one I recorded''.

\begin{tip}
Pick proteins with the same make-up as your recorded one: the same number of
chains, and the same kinds of extra parts (waters, ions, ligands). \pdb{1UBQ}
and \pdb{1AKI} both have one chain and crystal waters, and both worked with
the same template.
\end{tip}

\subsection{Stopping or carrying on after a failure}

Without \texttt{--batch-continue}, ProPrep stops at the first failure. The
proteins after it are marked \prompt{Not attempted} and the summary tells you
what to do:

\begin{screen}
Batch stopped after run 1 of 3. Pass --batch-continue to work through
failures instead of stopping at the first one.
\end{screen}

Use the default when a failure probably means your template is wrong, and
\texttt{--batch-continue} when you expect a few proteins to be awkward and
want the rest finished.

\section{Browser-recorded and terminal-recorded templates}

ProPrep's screens are the same in the browser (\texttt{proprep-web}) and in
the Terminal, but not quite: in the browser the structure viewer is docked
beside the terminal, so the PDB Filter does not ask whether you want to open
a viewer. A recording from one mode therefore has one question more or fewer
than the other.

ProPrep handles this for you. Each recording remembers where it was made, and
a batch is replayed in the same mode. You see which on the
\prompt{Replay mode:} line, either \prompt{plain terminal} or
\prompt{proprep-web (browser)}.

\begin{watchout}
Do not force the other mode with \texttt{--web} or \texttt{--no-web} unless
you know the recording suits it. Running a terminal-recorded template with
\texttt{--web} failed for us at exactly the question that differs:
\begin{screen}
Replay diverged: no recorded answer for prompt "Select chains ..."; the
recording expects a confirm here: 'Would you like to view the structure in
the interactive 3D viewer?'
\end{screen}
With \texttt{--no-web} the same template succeeded.
\end{watchout}

\section{Options at a glance}
\label{sec:batch-list}

\begin{tabularx}{\linewidth}{@{}lL@{}}
\toprule
\textbf{Option} & \textbf{What it does} \\
\midrule
\texttt{--create-template S.json} & Open a recorded session and mark blanks in
  it. \\
\texttt{--template-info T.json} & Print what a template contains: when it was
  made, its blanks and how many answers it holds. \\
\texttt{--validate-template T.json} & Check a template, and with
  \texttt{--batch-list} check that the list fills every blank. \\
\texttt{--batch-replay T.json} & Run a template. Needs
  \texttt{--batch-list}. \\
\texttt{--batch-list FILE} & The list: a CSV whose header row names the
  blanks, or a plain text file with one PDB code per line (which fills
  \texttt{input\_protein}). \\
\texttt{--output-dir DIR} & Where \file{run_01}, \file{run_02} and the
  reports go. Must already exist. \\
\texttt{--batch-continue} & Carry on after a protein fails. \\
\texttt{--web} / \texttt{--no-web} & Force browser or terminal mode. The
  default, taken from the recording, is almost always right. \\
\bottomrule
\end{tabularx}

\medskip
Each of these options warns you when it has nothing to act on, for example
\prompt{Warning: --batch-list has no effect without --batch-replay or
--validate-template}.

\section{If something goes wrong}

\begin{tabularx}{\linewidth}{@{}lL@{}}
\toprule
\textbf{What you see} & \textbf{What to do} \\
\midrule
\prompt{Setup error: [Errno 2] No such file or directory} for every run &
  Create the output folder first with \texttt{mkdir}. \\
\prompt{Error: --batch-replay requires --batch-list} & Add the list. \\
\prompt{File is not a template: ...} & You passed a plain session file.
  Make a template from it with \texttt{--create-template}. \\
\prompt{Variable 'water\_option' ...: no run supplies a value} & Your list
  has no column for that blank, or a row is empty. Nothing is run. \\
\prompt{No proteins found in input list: empty.txt} & The list has only
  comments or nothing at all. \\
\prompt{Replay diverged ...} for one protein & That protein differs from the
  recorded one (Section~\ref{sec:batch-choose}). Prepare it by hand, or
  record a template that fits it. \\
A run ``succeeded'' but the output is wrong & Check the values in your list; a
  value that is not a valid menu choice is accepted without a warning. \\
\bottomrule
\end{tabularx}

\begin{watchout}
A plain text list of UniProt or AlphaFold codes is read wrongly in version
1.23.0. A first line that starts with a letter is taken for a header row, so a
file beginning with \typed{P69905} gives \prompt{! Input column 'P69905'
matches no template variable} and then
\prompt{Variable 'input\_protein' ...: no run supplies a value}. PDB codes
like \typed{1UBQ} start with a digit and are fine. The workaround is to put
the line \typed{input\_protein} at the top of the file, which turns it into a
proper one-column table.
\end{watchout}

\begin{watchout}
When a protein fails, its Python error report is printed straight onto the
progress display, which looks alarming and scrolls the summary away. The same
text is saved in that run's \file{error.log}; read it there instead.
\end{watchout}
